On prépare une solution aqueuse en dissolvant une masse $m=\qty{271}{\mg}$ de
chlorure de fer(III) $\ce{FeCl3}$ anhydre dans un volume $V=\qty{250}{\mL}$
d'eau distillée.
\begin{enumerate}
    \item Écrire l'équation de dissolution de $\ce{FeCl3}$ dans
          l'eau.~\textbf{(0.5pt)}
    \item Déterminer la concentration $C$ de la solution obtenue en
          $\unit{\mol\per\L}$ puis en $\unit{\mol\per\cubic\m}$.~\textbf{(1pt)}

          \textbf{On donne~:} $M(\ce{FeCl3})=\qty{162.5}{\g\per\mol}$
    \item Déterminer les concentrations effectives des espèces chimiques qui se
          trouvent dans la solution.~\textbf{(1.5pt)}
    \item En déduire la conductivité $\sigma$ de la solution.~\textbf{(1pt)}

          \textbf{On donne~:} les conductivités molaires ioniques~: \\
          $\lambda(\ce{Cl^-})=\qty{7.63e-3}{\lambda}$ et
          $\lambda(\ce{Fe^3+})=\qty{2.04e-2}{\lambda}$.
\end{enumerate}
